\noindent\underline{证明.}
鉴于 $b_j$ 的定义 (5.4.3), 只需证明
\begin{equation}\tag{5.4.6}
\left\{
\begin{aligned}
a_0&=1,\\
\sum_{i=0}^{a}(-1)^iS_i\,a_{j-i}&=0
&&\text{若 }j>0.
\end{aligned}
\right.
\end{equation}
\EditorialFootnote{下标约定: (5.4.6) 中出现负下标时, 取 $a_q=0$（$q<0$）. 这正是从 (5.4.1) 取 $T^{a-1}$ 的系数时, 次数小于 $a-1$ 的幂所给出的零系数; 原文未另行写出这一约定.}
显然 $a_0=1$. 对固定的 $j>0$, 利用 (5.4.1), 在关系式
\begin{equation}\tag{5.4.7}
\sum_{i=0}^{a}(-1)^iS_iT^{a-1+j-i}=0
\qquad\text{在 }B.
\end{equation}
中取 $T^{a-1}$ 的系数即得第二个关系式.
\footnote{印本在 (5.4.7) 中漏去了求和上界 $a$; 此关系式为 (5.4.2) 乘以 $T^{j-1}$.}
此关系式由 (5.4.2) 乘以 $T^{j-1}$ 得到, 引理得证.

\subsubsection*{推论 5.5.}
\underline{有}
\begin{equation}\tag{5.5.1}
\deg(\check X)
=\frac{(-1)^n}{\deg(\varphi)}
\int_X\frac{(1+H)^{r-1}}{C(N)}.
\end{equation}

\noindent\underline{证明.}
取 $M=N$ 为 $X\hookrightarrow\mathbb P^r$ 的法层, 对 (5.2.8) 应用 (5.3.3).

对光滑拟射影簇 $S$, 将其全 Chern 类记作
\begin{equation}\tag{5.5.2}
C(S)\stackrel{\mathrm{def}}{=}C(T_S)
\end{equation}
其中 $T_S$ 为 $S$ 的切丛.

将正合列 (5.1.2) 对偶化, 得到
\begin{equation}\tag{5.5.3}
C(\mathbb P^r)=(1+H)^{r+1}
\qquad\text{在 }\operatorname{CH}(\mathbb P^r).
\end{equation}
将正合列 (5.1.1) 对偶化, 得到
\begin{equation}\tag{5.5.4}
C(N)\,C(X)=C(\mathbb P^r)|_X
\qquad\text{在 }\operatorname{CH}(X).
\end{equation}
结合 (5.5.3) 和 (5.5.4), 公式 (5.5.1) 化为:

\subsubsection*{推论 5.6.}
\underline{有}
\begin{equation}\tag{5.6.1}
\deg(\check X)
=\frac{(-1)^n}{\deg(\varphi)}
\int_X\frac{C(X)}{(1+H)^2}.
\end{equation}
\footnote{印本写为 $(-1)^r$. 将 (5.5.3) 和 (5.5.4) 代入 (5.5.1) 并不改变其符号 $(-1)^n$; 紧接在后的公式 (5.7.3) 也确认相同的指数.}

\subsubsection*{5.7.}
对光滑射影簇 $S$, 定义其 \underline{Euler--Poincaré 示性数}为
\begin{equation}\tag{5.7.1}
\chi(S)=\int_SC(S).
\end{equation}

\subsubsection*{命题 5.7.2.}
\underline{设 $Y$ 为 $X$ 的光滑超平面截面, $\Delta$ 为 $X$ 与 $\mathbb P^r$ 中余维二线性子空间的光滑截面. 则}
\begin{equation}\tag{5.7.3}
\deg(\check X)
=\frac{(-1)^n}{\deg(\varphi)}
\bigl[\chi(X)+\chi(\Delta)-2\chi(Y)\bigr].
\end{equation}

